\chapter{Materials and Methods}\label{chap:materials-methods}

\section{Research Design}\label{sec:methods-research-design}

This thesis is organized around a reproducible experimental protocol. Each experiment is described by its data-generating source, preprocessing transformation, feature representation, model family, evaluation rule, and recorded outcomes.

% !TEX root = main.tex
\subsection{Dataset preprocessing and construction}

Let $\mathcal{S}=\{\mathrm{Bhandari},\mathrm{CDBProm},\mathrm{CNNProm},\mathrm{GraphPro}\}$ denote the set of source studies considered in this thesis. For each source $s\in\mathcal{S}$, the experimental data are represented as a finite labeled set $\mathcal{D}_{s}=\{(x_i,y_i)\}_{i=1}^{n_s}\subset\mathcal{X}_L\times\mathcal{Y}$, with the input space $\mathcal{X}_L$ and label space $\mathcal{Y}$ of Section~\ref{sec:bg-formulation}. After preprocessing, every sequence window of a dataset has the same length $L$, fixed by its biological domain (81 bp for prokaryotes and 251 bp for eukaryotes), so the experimental unit is one fixed-length sequence-level example.

The four source studies differ in organism coverage, sequence-window definition,
and negative-set construction. In particular, CDBProm, CNNProm, and GraphPro
provide both labeled classes, whereas the downloaded Bhandari collections contain
promoter sequences only. Bhandari experiments therefore require a separately
declared synthetic negative-generation rule. Table~\ref{tab:active-source-provenance}
summarizes the study-level provenance used to define the source set $\mathcal{S}$
before describing the transformations applied to obtain the standardized labeled
sets $\mathcal{D}_{s}$.

\begin{table}[htbp]
\centering
\small
\caption{Study-level provenance of the source datasets used in this thesis. The table identifies the publications or resources from which the promoter and non-promoter sequence sets were derived.}
\label{tab:active-source-provenance}
\begin{tabular}{p{0.18\textwidth}p{0.42\textwidth}p{0.30\textwidth}}
\hline
Source & Publication or resource & Recorded origin \\
\hline
Bhandari et al. & Comparison of machine learning and deep learning techniques in promoter prediction across diverse species & Promoter-only public data resource associated with the study \\
CDBProm & CDBProm: the Comprehensive Directory of Bacterial Promoters, Martinez et al., 2024 & CDBProm website and training-sequence archive \\
CNNProm & Recognition of Prokaryotic and Eukaryotic Promoters Using Convolutional Deep Learning Neural Networks, Umarov and Solovyev, 2016 & CNNPromoterData resource and Softberry website \\
GraphPro & GraphPro: An interpretable graph neural network-based model for identifying promoters in multiple species, Zhang et al., 2024 & GraphPro public dataset resource \\
\hline
\end{tabular}
\end{table}

The source descriptions also provide methodological context that affects interpretation. For example, CNNProm reports bacterial non-promoter sequences sampled from protein-coding genes on the opposite strand, whereas its promoter sets are linked to resources such as RegulonDB, DBTBS, EPDnew, and DBTSS. GraphPro includes human and mouse eukaryotic datasets and bacterial datasets for six organisms. These differences motivate source-aware evaluation, since the data-generating distribution $p_{\text{data},s}$ may differ across studies.

Source labels were audited separately from sequence-level integrity.
Deduplication can remove an observable contradiction, but it cannot establish
that a file named as a negative class contains biologically defensible
non-promoters. In the distributed GraphPro bacterial files, the class is
perfectly determined by the strand symbol in five of six organism folders:
every positive header denotes the plus strand and every negative header denotes
the minus strand. Because promoters can occur on either strand, this pattern
does not support interpreting strand as promoter status. GraphPro
\textit{E. coli} additionally contains 1,345 records with identical sequence,
accession, coordinate, strand, and complete FASTA header in both classes.

The audit also corrected an organism-mapping error in the CapsProm ingestion
code. The GraphPro folder \texttt{Cpne} had been named as
\textit{Clostridium phytofermentans}, but its records use accession
\texttt{NC\_000922.1} and explicitly identify \textit{Chlamydophila pneumoniae}
CWL029. Future reconstruction uses the header-supported organism identity;
legacy v3 artifacts retain their original identifier only for provenance.

Consequently, exact-sequence cleanliness and source-label validity are treated
as distinct eligibility gates. The confirmatory scope contains the 11 CDBProm
and seven CNNProm datasets. All eight GraphPro datasets are retained as
exploratory provenance until their negative-set construction can be verified or
independently reconstructed. No GraphPro record is automatically relabeled.

The treatment path preserves the released GraphPro files as an exploratory
reproduction benchmark and defines a distinct derived family,
\texttt{graphpro-reconstructed-v1}. For bacterial organisms, positives must be
recreated from independently documented experimental TSS coordinates and the
versioned reference assembly. For human and mouse, every \texttt{FP}
identifier must first be resolved against a dated EPD release and its genome
assembly. Because EPD is a promoter/TSS resource, membership in the released
GraphPro \texttt{negative} directory is not accepted as evidence of a
non-promoter label.

The identifier-resolution stage uses the November 2023 EPDnew archive. Human
identifiers are resolved against HsNC EPDnew version 001 on GRCh38/hg38 and
mouse identifiers against MmNC EPDnew version 001 on GRCm38/mm10. Coding EPDnew
collections are retained only as alternative annotations. The released file
class, the \texttt{FP} prefix, and the embedded header label remain provenance
fields and are not used as independent biological evidence. Coordinate
resolution precedes sequence recreation and negative-control generation, which
remain separate eligibility gates.

Complete sequence recreation subsequently established the released eukaryotic
construction rule. Relative to the strand-oriented EPD TSS, positive windows
span $[-80,1)$ and negative windows span $[1000,1081)$. All 8,370 released
Human and Mouse records match the frozen hg38/mm10 assemblies exactly. The
downstream windows are therefore treated as a reproducible, source-defined
control condition, but not as experimentally verified non-promoters or as the
primary GC-matched control policy.

For the historical downstream-control condition, treatment is applied after
the released-data exact-integrity stage. Strand-invariant cross-label
collisions are removed bilaterally and within-class sequence/reverse-complement
duplicates retain only their first source occurrence. A negative window is
excluded when its interval, expanded by 100 nt on each side, contains any TSS
from the union of the coding and non-coding EPDnew collections. Only the invalid
control is removed because the Mouse source does not define a complete
one-to-one positive--control pairing. No compensatory downsampling is applied;
records sharing an EPD identifier are instead preserved as an indivisible group
for future partitioning.

The version-2 historical condition adds an external-positive control check. A
negative control is unsuitable when its exact sequence or reverse complement
is labeled positive in a hash-bound dataset in the active comparison scope.
The control alone is removed under this rule, while the external dataset and
its label remain unchanged. The exclusion ledger records both dataset
identities, the external label, match orientation, source index, identifier,
and sequence hash. This policy prevents an explicit control contradiction but
does not establish which source annotation is biologically correct.

Each verified positive is represented by the strand-oriented interval from 60
nt upstream through 20 nt downstream of the TSS. Its primary control is sampled
deterministically from the same assembly, with the same length and matched
strand and GC distributions, outside a frozen exclusion region around every
known TSS. Candidate controls that overlap an exclusion interval, duplicate a
positive or another control, contain ambiguous symbols, or violate the frozen
GC tolerance are rejected and counted. Positive--control pairs and records
sharing one TSS or source experiment form indivisible partition groups.

For the Human/Mouse primary candidate, matching is defined at the pair level:
each control is sampled from the same chromosome, oriented on the same strand,
and differs by at most two G/C bases from its 81-nt promoter. Exact equality
first failed the frozen 100,000-candidate ceiling for one Human promoter; a
separate one-base relaxation then failed for one Mouse promoter. The two-base
random probe failed for that same Mouse promoter, whose window contains 75 G/C
bases (92.59\%). These feasibility decisions used source composition only and
did not inspect model outcomes.

The final candidate therefore retains the two-base tolerance and invokes a
bounded-memory exhaustive chromosome scan only after random exhaustion.
Eligible fallback windows are chosen by deterministic reservoir sampling while
all TSS-margin, sequence-identity, ambiguous-DNA, and control-overlap exclusions
remain active. This procedure produced 2,332 Human pairs and 2,571 Mouse pairs
after bilateral removal of seven and five duplicate-positive pairs,
respectively. Two Mouse promoters required fallback:
\texttt{0610009L18Rik\_1} (75/81 G/C, paired with 74/81) and
\texttt{Gm15327\_1} (76/81, paired with 76/81). Random and exhaustive search
effort are recorded in separate ledger fields.
Here, pairing denotes a construction and grouping relation rather than
biological homology. The positive and control share a stable pair identifier,
are filtered together, and must be assigned to the same future fold. Controls
also cannot match an active external positive in either orientation. Exact
pair-level GC matching strongly constrains the univariate GC-count difference,
but does not establish equivalence in aggregate, higher-order, or positional
composition; these properties are audited separately before eligibility.

In the descriptive audit, the mean positive-minus-negative GC-fraction
difference was 0.00289 for Human and 0.00405 for Mouse, compared with
approximately 0.11 in the historical downstream-control condition. Neither the
two-sided Mann--Whitney nor Kolmogorov--Smirnov comparison indicated an
aggregate GC-distribution difference after Holm correction. This result
supports reduction of the intended univariate shortcut, but it does not imply
biological equivalence: the largest positional nucleotide-composition
difference occurs at the TSS, and dinucleotide, chromatin, gene-proximity, and
unobserved-promoter effects remain possible. Exact dependency analysis found
only same-label positive sharing with CNNProm and between the Human and Mouse
candidates; no opposite-label identity was observed. These dependencies must
be respected by grouping or by a predeclared cross-source separation policy.

Within each candidate, the strand-invariant Hamming-distance audit found 53
Human pairs and 49 Mouse pairs with a neighbor at distance at most eight. The
largest connected component contained six pairs in either organism. Human had
three cross-label near-neighbor edges at this threshold, including one at
distance four, whereas Mouse had none. These relations are not exact label
collisions; they motivate assigning the entire connected pair component to one
fold and auditing train--test separation after partitioning.

A strand-invariant near-duplicate graph additionally joins pairs whenever any of
their sequences differ by at most eight substitutions in either orientation.
Each connected component is assigned wholly to training or testing by a
deterministic component hash (seed 42), targeting an 80/20 division. The
partition is accepted only after a direct train--test audit confirms no exact,
reverse-complement, or distance-eight neighbor. A separate CDS-derived negative
condition is reserved as a sensitivity analysis,
because the definition of the negative class changes the task itself.

Reinstatement is evaluated per organism rather than collectively. It requires
verified source coordinates and assemblies, immutable hashes, a complete
candidate-rejection and filtering ledger, group-safe partitions, exact and
near-duplicate audits, distribution diagnostics, and manual verification of a
record sample. Only a new eligibility and dependency manifest may admit a
reconstructed dataset to confirmatory experiments; predictive performance is
never an eligibility criterion.

Because the source studies provide heterogeneous sequence collections, all records were first mapped into a common labeled representation. Promoter sequences were assigned label $1$, whereas non-promoter sequences were assigned label $0$. For each source $s$, this step defines an initial labeled collection $\mathcal{D}^{0}_{s}=\{(\tilde x_i,y_i)\}_{i=1}^{\tilde n_s}$, where $\tilde x_i$ denotes the original sequence window before length normalization.

A deterministic cropping function $\operatorname{crop}_s$ was then applied to each original sequence $\tilde x_i$ to obtain a standardized sequence $x_i=\operatorname{crop}_s(\tilde x_i)$ with $|x_i|=81$. The target window is defined relative to a source-specific transcription start site (TSS) coordinate and contains 60 bp upstream of the TSS, the TSS nucleotide itself, and 20 bp downstream. This window definition was adopted to maintain compatibility with the promoter-region representation used by the GraphPro and CNNProm studies for prokaryotic datasets, while enforcing a common input length across all sources. Because the original sequence windows differ across studies, the TSS coordinate used by $c_s$ was defined per source and original window length.

Table~\ref{tab:tss-cropping-parameters} summarizes the source-specific TSS
coordinates used to obtain the standardized 81 bp windows.

\begin{table}[htbp]
\centering
\caption{Source-specific parameters of the deterministic cropping function $\operatorname{crop}_s$ used to map promoter and non-promoter sequences to 81 bp windows. The final window contains 60 bp upstream of the transcription start site, the TSS nucleotide, and 20 bp downstream.}
\label{tab:tss-cropping-parameters}
\begin{tabular}{lrrr}
\hline
Source & Original length & TSS position & Final length \\
\hline
GraphPro & 81 & 60 & 81 \\
CNNProm & 81 & 60 & 81 \\
CNNProm & 251 & 200 & 81 \\
CDBProm & 100 & 60 & 81 \\
CDBProm & 1001 & 500 & 81 \\
Bhandari et al. & 1000 & 700 & 81 \\
\hline
\end{tabular}
\end{table}

After cropping, promoter and non-promoter examples from CDBProm, CNNProm, and
GraphPro were paired within each source-defined dataset. Each resulting subset
is represented as $\mathcal{D}_{s,d}=\mathcal{P}_{s,d}\cup\mathcal{N}_{s,d}$,
where $\mathcal{P}_{s,d}=\{(x_i,1)\}$ and
$\mathcal{N}_{s,d}=\{(x_i,0)\}$ denote the positive and negative examples for
dataset $d$ from source $s$. Their corrected collections receive deterministic
train/test assignments, which are then subject to integrity auditing and must
be replaced by exact-sequence group assignments whenever row-level assignment
produces partition overlap.

\subsubsection{Exact-sequence identity and correction policies}

All exact-sequence checks use the uppercase nucleotide string as the identity
key after removal of FASTA line breaks and surrounding whitespace. Four cases
are distinguished. A \emph{same-label duplicate} is a repeated identity within
one class; a \emph{cross-label collision} is an identity associated with both
labels; a \emph{partition overlap} is an identity found in both training and
test; and a \emph{generation duplicate} is an artificial repetition introduced
by a preprocessing execution, such as appending to an existing output file.
These categories are reported separately because they have different causes
and require different corrections.

Generation duplicates are prevented by overwrite-only output and an
idempotence test. Cross-label collisions are removed bilaterally as described
below. Same-label multiplicity is then removed by retaining one occurrence of
each normalized sequence. After strict-DNA filtering, partitioning is performed
over unique exact-sequence groups, rather than source rows, so that no identity
can occur in both training and test. A source-provided fixed split that violates
either rule is retained only for provenance and is ineligible for confirmatory
evaluation until rebuilt.

Cropping can map distinct source records to the same 81 bp sequence. Therefore,
before partitioning, the pipeline identifies the collision set
$\mathcal{C}_{s,d}=\{x:(x,0)\in\mathcal{D}_{s,d}\land
(x,1)\in\mathcal{D}_{s,d}\}$. Every occurrence of a sequence in
$\mathcal{C}_{s,d}$ is removed from both classes. This bilateral policy avoids
arbitrary relabeling and prevents an identical input from representing opposing
targets. A filtering manifest records, for every dataset, the row counts before
and after cropping, collision removal, partitioning, and strict-alphabet
filtering, including the number removed at each stage.

The Bhandari source follows a distinct control workflow because no genomic
non-promoter collection was supplied. Exact promoter windows are deduplicated
before partitioning. For each retained promoter, one deterministic
mononucleotide shuffle is generated, preserving the counts of A, C, G, and T
while changing their order. A promoter and its derived shuffle form an
indivisible group, and groups are assigned to train or test with seed 42. This
construction prevents exact train--test overlap and preserves a one-to-one
class balance, but it defines a shuffle-control task rather than a biological
promoter-versus-genomic-non-promoter benchmark.

Table~\ref{tab:sequence-integrity-policies} fixes the identification and
resolution rule for each exact-sequence integrity condition. These rules apply
before any feature representation is fitted.

\begin{table}[htbp]
\centering
\caption{Exact-sequence integrity policies applied during dataset construction. Identification uses normalized sequence identity; counts and removals are recorded per class and processing stage.}
\label{tab:sequence-integrity-policies}
\begin{tabular}{p{0.20\textwidth}p{0.31\textwidth}p{0.39\textwidth}}
\toprule
Condition & Identification & Resolution policy \\
\midrule
Generation duplicate & Output grows when the same deterministic step is rerun & Overwrite outputs and verify idempotence \\
Same-label duplicate & Same sequence occurs more than once with one label & Retain one normalized sequence before partitioning; preserve source counts in provenance records \\
Cross-label collision & Same sequence occurs with labels 0 and 1 & Remove every occurrence from both classes before partitioning \\
Partition overlap & Same sequence occurs in training and test & Rebuild using indivisible sequence groups; reject the fixed split meanwhile \\
Invalid DNA & Length differs from 81 or symbols fall outside $\Sigma_{\mathrm{DNA}}$ & Remove before partition assignment and record the affected class \\
\bottomrule
\end{tabular}
\end{table}

\subsubsection{Quality control and final dataset definition}

An alphabet filter was then applied before partitioning to define the final
model-ready sequence pool. Let $\Sigma_{\mathrm{DNA}}=\{\mathrm{A},\mathrm{C},\mathrm{G},\mathrm{T}\}$ be the accepted DNA
alphabet. The filtering predicate $\operatorname{keep}(x_i,y_i)$ is true, and the example
is retained, if and only if $x_i\in\Sigma_{\mathrm{DNA}}^{81}$ and
$y_i\in\mathcal{Y}$. This policy ensures that all
subsequent encodings, including one-hot, k-mer, and physicochemical encodings,
operate over the same unambiguous DNA alphabet. In the canonical CDBProm,
CNNProm, and GraphPro build, this stage removed eight rows from five datasets:
six positive and two negative examples. The earlier legacy inventory removed
338 examples from 12 subsets and is retained only for provenance. The active
stage counts are reported in Table~\ref{tab:results-filtering-stage-counts}, and
the complete per-dataset ledger is reported in
Table~\ref{tab:appendix-filtering-stage-counts}.

The complete preprocessing workflow is summarized in Figure~\ref{fig:within-source-preprocessing-workflow}, which shows how heterogeneous source collections were transformed into validated, model-ready labeled sequence sets.

\begin{figure}[htbp]
\centering
\begin{tikzpicture}[
    node distance=0.85cm,
    process/.style={rectangle, rounded corners, draw=black, align=center, minimum width=0.78\textwidth, minimum height=0.82cm, font=\small},
    data/.style={trapezium, trapezium left angle=70, trapezium right angle=110, draw=black, align=center, minimum width=0.74\textwidth, minimum height=0.82cm, font=\small},
    output/.style={rectangle, draw=black, align=center, minimum width=0.78\textwidth, minimum height=0.82cm, font=\small},
    arrow/.style={-{Latex[length=2.4mm]}, thick}
]
\node[data] (sources) {Source studies $\mathcal{S}$\\Bhandari et al., CDBProm, CNNProm, and GraphPro};
\node[process, below=of sources] (fasta) {Source normalization\\Map heterogeneous records to labeled sequence sets};
\node[process, below=of fasta] (crop) {TSS-centered cropping\\Extract 81 bp windows: 60 bp upstream, TSS base, and 20 bp downstream};
\node[process, below=of crop] (pair) {Class construction\\Use source negatives or declared Bhandari shuffles};
\node[process, below=of pair] (collisions) {Cross-label collision filtering\\Remove every conflicting sequence from both classes};
\node[process, below=of collisions] (deduplicate) {Within-class deduplication\\Retain one row per normalized sequence and label};
\node[process, below=of deduplicate] (clean) {Alphabet filtering\\Retain only $x_i\in\Sigma_{\mathrm{DNA}}^{81}$};
\node[process, below=of clean] (csv) {Partition definition\\Assign deterministic splits; keep exact-sequence groups indivisible};
\node[process, below=of csv] (validate) {Constraint validation\\Check length, labels, class counts, duplicates, and overlap};
\node[output, below=of validate] (ready) {Audited labeled partitions\\Eligible sets proceed to feature encoding and experiments};

\draw[arrow] (sources) -- (fasta);
\draw[arrow] (fasta) -- (crop);
\draw[arrow] (crop) -- (pair);
\draw[arrow] (pair) -- (collisions);
\draw[arrow] (collisions) -- (deduplicate);
\draw[arrow] (deduplicate) -- (clean);
\draw[arrow] (clean) -- (csv);
\draw[arrow] (csv) -- (validate);
\draw[arrow] (validate) -- (ready);
\end{tikzpicture}
\caption{Preprocessing workflow used to construct the promoter prediction datasets. The diagram distinguishes source-provided non-promoters from the declared Bhandari shuffle-control path, then summarizes cropping to 81 bp, mandatory bilateral removal of cross-label collisions, within-class deduplication, strict-alphabet filtering, deterministic group-safe partitioning, and validation before feature encoding and model training. Each filtering stage records its input, removed, and remaining row counts in the filtering manifest.}
\label{fig:within-source-preprocessing-workflow}
\end{figure}


The legacy strict-DNA inventory is summarized in
Table~\ref{tab:within-source-final-summary}. Its Bhandari counts are retained for
provenance only because they include replicated promoters and legacy synthetic
negatives. The accepted Bhandari human shuffle-control is reported separately
in Table~\ref{tab:bhandari-leakage-controlled-counts}. The remaining
study-specific tables provide the source-defined train/test counts for CDBProm
in Table~\ref{tab:cdbprom-within-source-final}, CNNProm in
Table~\ref{tab:cnnprom-within-source-final}, and GraphPro in
Table~\ref{tab:graphpro-within-source-final}.

\begin{table}[htbp]
\centering
\caption{Legacy strict-DNA inventory after alphabet cleaning. The Bhandari row contains legacy replicated positives and synthetic negatives and is retained only for preprocessing provenance; it is not the active Bhandari experimental dataset.}
\label{tab:within-source-final-summary}
\begin{tabular}{lrrrr}
\hline
Study & Subsets & Final $n^{+}$ & Final $n^{-}$ & Final total \\
\hline
Bhandari et al. & 6 & 583819 & 583819 & 1167638 \\
CDBProm & 22 & 46959 & 62420 & 109379 \\
CNNProm & 14 & 141948 & 297680 & 439628 \\
GraphPro & 16 & 47799 & 64480 & 112279 \\
\hline
\end{tabular}
\end{table}

Class balance in the legacy inventory is summarized in
Table~\ref{tab:within-source-class-balance}. The apparent 50/50 Bhandari balance
is a consequence of synthetic negative construction rather than a property of
the downloaded source. CDBProm, CNNProm, and GraphPro retain their
source-provided class distributions.

\begin{table}[htbp]
\centering
\caption{Class balance of the legacy strict-DNA inventory by study. Percentages are computed from the cleaned sequence counts; the Bhandari balance reflects synthetic negative construction, whereas the other sources include source-provided non-promoters.}
\label{tab:within-source-class-balance}
\begin{tabular}{lrrrr}
\hline
Study & Promoter (\%) & Non-promoter (\%) & $n^{+}/n^{-}$ ratio & Final total \\
\hline
Bhandari et al. & 50.0 & 50.0 & 1.00 & 1167638 \\
CDBProm & 42.9 & 57.1 & 0.75 & 109379 \\
CNNProm & 32.3 & 67.7 & 0.48 & 439628 \\
GraphPro & 42.6 & 57.4 & 0.74 & 112279 \\
\hline
\end{tabular}
\end{table}

Table~\ref{tab:bhandari-within-source-final} documents the legacy Bhandari
partitions for provenance. These files are excluded from current model
comparisons because exact duplicates occur within partitions and across the
train--test boundary.

\begin{table}[htbp]
\centering
\small
\caption{Legacy unambiguous-DNA Bhandari datasets retained for provenance. The rows include synthetic negatives and must not be used as leakage-controlled biological benchmarks.}
\label{tab:bhandari-within-source-final}
\begin{tabular}{p{0.40\textwidth}lrrr}
\hline
Dataset & Split & Final $n^{+}$ & Final $n^{-}$ & Final total \\
\hline
eukaryotic Human\_promoter & test & 49188 & 49187 & 98375 \\
eukaryotic Human\_promoter & train & 196756 & 196757 & 393513 \\
eukaryotic Plant\_promoters & test & 63900 & 63901 & 127801 \\
eukaryotic Plant\_promoters & train & 255609 & 255608 & 511217 \\
eukaryotic Yeast\_promoter & test & 3673 & 3674 & 7347 \\
eukaryotic Yeast\_promoter & train & 14693 & 14692 & 29385 \\
\hline
\end{tabular}
\end{table}

The corrective human workflow starts from 61,546 raw promoter records. It
retains 61,516 records of the expected 1,000 bp length, removes 30 cropped
windows containing non-strict symbols, and collapses 27,561 repeated promoter
occurrences. The resulting 33,925 unique promoters each receive one valid
shuffle control. Table~\ref{tab:bhandari-leakage-controlled-counts} reports the
accepted group-level split. Every sequence is unique, no label conflict is
present, and no exact sequence is shared between train and test.

\begin{table}[htbp]
\centering
\caption{Accepted leakage-controlled Bhandari human shuffle-control dataset. Each promoter and its deterministic mononucleotide shuffle remain in the same partition; the negative class is synthetic and does not represent genomic non-promoters.}
\label{tab:bhandari-leakage-controlled-counts}
\begin{tabular}{lrrrr}
\hline
Partition & Promoter & Synthetic negative & Total & Unique sequences \\
\hline
Train & 27140 & 27140 & 54280 & 54280 \\
Test & 6785 & 6785 & 13570 & 13570 \\
Total & 33925 & 33925 & 67850 & 67850 \\
\hline
\end{tabular}
\end{table}

Table~\ref{tab:cdbprom-within-source-final} reports the final CDBProm train/test partitions for each prokaryotic dataset.

\begin{table}[htbp]
\centering
\small
\caption{Final unambiguous-DNA CDBProm within-source datasets after preprocessing and cleaning.}
\label{tab:cdbprom-within-source-final}
\begin{tabular}{p{0.40\textwidth}lrrr}
\hline
Dataset & Split & Final $n^{+}$ & Final $n^{-}$ & Final total \\
\hline
prokaryotic anabaena\_sp & test & 1510 & 2009 & 3519 \\
prokaryotic anabaena\_sp & train & 6041 & 8035 & 14076 \\
prokaryotic bacillus\_amyloliquefaciens & test & 636 & 845 & 1481 \\
prokaryotic bacillus\_amyloliquefaciens & train & 2541 & 3379 & 5920 \\
prokaryotic chlamydophila\_pneumoniae & test & 243 & 323 & 566 \\
prokaryotic chlamydophila\_pneumoniae & train & 969 & 1293 & 2262 \\
prokaryotic corynebacterium\_glutamicum & test & 345 & 458 & 803 \\
prokaryotic corynebacterium\_glutamicum & train & 1380 & 1830 & 3210 \\
prokaryotic helicobacter\_pylori & test & 429 & 568 & 997 \\
prokaryotic helicobacter\_pylori & train & 1707 & 2272 & 3979 \\
prokaryotic klebisiella\_pneumoniae & test & 798 & 1061 & 1859 \\
prokaryotic klebisiella\_pneumoniae & train & 3189 & 4243 & 7432 \\
prokaryotic mycrobacterium\_tuberculosis & test & 1063 & 1414 & 2477 \\
prokaryotic mycrobacterium\_tuberculosis & train & 4253 & 5654 & 9907 \\
prokaryotic pseudonomas\_aeruginosa & test & 1232 & 1640 & 2872 \\
prokaryotic pseudonomas\_aeruginosa & train & 4930 & 6556 & 11486 \\
prokaryotic salmonella\_typhimurium & test & 589 & 779 & 1368 \\
prokaryotic salmonella\_typhimurium & train & 2354 & 3117 & 5471 \\
prokaryotic streptomyces\_coelicolor & test & 1567 & 2085 & 3652 \\
prokaryotic streptomyces\_coelicolor & train & 6266 & 8339 & 14605 \\
prokaryotic synechocistis & test & 984 & 1304 & 2288 \\
prokaryotic synechocistis & train & 3933 & 5216 & 9149 \\
\hline
\end{tabular}
\end{table}

Table~\ref{tab:cnnprom-within-source-final} reports the final CNNProm train/test partitions across eukaryotic and prokaryotic subsets.

\begin{table}[htbp]
\centering
\small
\caption{Final unambiguous-DNA CNNProm within-source datasets after preprocessing and cleaning.}
\label{tab:cnnprom-within-source-final}
\begin{tabular}{p{0.40\textwidth}lrrr}
\hline
Dataset & Split & Final $n^{+}$ & Final $n^{-}$ & Final total \\
\hline
eukaryotic Arabidopsis\_non\_tata & test & 3543 & 9167 & 12710 \\
eukaryotic Arabidopsis\_non\_tata & train & 14172 & 36665 & 50837 \\
eukaryotic Arabidopsis\_tata & test & 898 & 2304 & 3202 \\
eukaryotic Arabidopsis\_tata & train & 3593 & 9212 & 12805 \\
eukaryotic Human\_non\_tata & test & 11887 & 22185 & 34072 \\
eukaryotic Human\_non\_tata & train & 47543 & 88739 & 136282 \\
eukaryotic Mouse\_non\_tata & test & 9770 & 19858 & 29628 \\
eukaryotic Mouse\_non\_tata & train & 39079 & 79430 & 118509 \\
eukaryotic Mouse\_tata & test & 1566 & 2824 & 4390 \\
eukaryotic Mouse\_tata & train & 6261 & 11296 & 17557 \\
prokaryotic Bacillus & test & 224 & 800 & 1024 \\
prokaryotic Bacillus & train & 895 & 3200 & 4095 \\
prokaryotic Ecoli & test & 504 & 2400 & 2904 \\
prokaryotic Ecoli & train & 2013 & 9600 & 11613 \\
\hline
\end{tabular}
\end{table}

Table~\ref{tab:graphpro-within-source-final} reports the final GraphPro train/test partitions across eukaryotic and prokaryotic subsets.

\begin{table}[htbp]
\centering
\small
\caption{Final unambiguous-DNA GraphPro within-source datasets after preprocessing and cleaning.}
\label{tab:graphpro-within-source-final}
\begin{tabular}{p{0.40\textwidth}lrrr}
\hline
Dataset & Split & Final $n^{+}$ & Final $n^{-}$ & Final total \\
\hline
eukaryotic Human & test & 1404 & 1871 & 3275 \\
eukaryotic Human & train & 5613 & 7485 & 13098 \\
eukaryotic Mouse & test & 1108 & 1477 & 2585 \\
eukaryotic Mouse & train & 4430 & 5907 & 10337 \\
prokaryotic bacillus\_amyloliquefaciens & test & 557 & 743 & 1300 \\
prokaryotic bacillus\_amyloliquefaciens & train & 2227 & 2969 & 5196 \\
prokaryotic campylobacter\_jejuni & test & 808 & 1077 & 1885 \\
prokaryotic campylobacter\_jejuni & train & 3230 & 4307 & 7537 \\
prokaryotic clostridium\_phytofermentans & test & 302 & 345 & 647 \\
prokaryotic clostridium\_phytofermentans & train & 1207 & 1379 & 2586 \\
prokaryotic escherichia\_coli & test & 1512 & 3099 & 4611 \\
prokaryotic escherichia\_coli & train & 6045 & 12397 & 18442 \\
prokaryotic leptospira\_interrogans & test & 1707 & 1693 & 3400 \\
prokaryotic leptospira\_interrogans & train & 6828 & 6767 & 13595 \\
prokaryotic rhodobacter\_capsulatus & test & 2164 & 2593 & 4757 \\
prokaryotic rhodobacter\_capsulatus & train & 8657 & 10371 & 19028 \\
\hline
\end{tabular}
\end{table}

These unambiguous-DNA sequence-level datasets provide the input for the feature
engineering stage. In the subsequent experimental protocol, each 81 bp sequence
can be represented using k-mer frequency vectors, one-hot matrices,
physicochemical property matrices, DNA-BERT embeddings, or controlled
combinations of these feature blocks depending on the classifier architecture.


\section{Reproducible Experimental Protocol}\label{sec:methods-protocol-overview}

The experimental protocol separates data validation, feature construction, model training, evaluation, and reporting. Each experiment is specified formally by the configuration tuple defined in Section~\ref{sec:protocol-configuration}, which separates the dataset, the representation, the model, its learned parameters, its hyperparameter configuration, the evaluation outcomes, and the random seed.

\section{Dataset Scope and Provenance}\label{sec:methods-dataset-scope}

The audited inventory contains source-defined datasets derived from Bhandari et
al., CDBProm, CNNProm, and GraphPro. The active confirmatory scope is restricted
to the 18 eligible CDBProm and CNNProm identities; GraphPro is retained for
exploratory provenance pending independent validation of its labels. Source provenance is reported
separately from feature reference tables to avoid mixing dataset origin with
encoding resources. CDBProm and CNNProm provide documented promoter and
non-promoter examples. Bhandari provides promoter sequences only and is treated
as a separate synthetic-negative control rather than as evidence of performance
against genomic non-promoters.

% !TEX root = main.tex
\subsection{Data acquisition and provenance}\label{sec:methods-data-acquisition}

A comparison between learning methods is only as interpretable as the data on
which it is made. Promoter benchmarks are distributed as files in code
repositories, journal supplements, and institutional servers. These copies can
change without notice, and several of them differ from the publications that
describe them. The data used in this thesis are therefore treated as an object
whose origin must be demonstrable, following the findability, accessibility,
interoperability and reusability principles for scientific data
\citep{wilkinson2016fair}. Three requirements guided the procedure: each
experimental example must be traceable to an identified upstream file; the
upstream copy must be identified independently of the moment at which it was
obtained; and any later acquisition must be distinguishable from the first.

\paragraph{Reference collection.}
The datasets released with the evaluated studies
\citep{umarov2017recognition,bhandari2021comparison,martinez2024cdbprom,zhang2024graphpro},
together with primary references used for reconstruction, such as the
Eukaryotic Promoter Database \citep{meylan2020epd}, were acquired again directly
from their distribution points into a single versioned reference collection.
Each source was identified by an access identifier that cannot drift. Code
repositories were pinned to a full commit identifier, and every file had to
reproduce the content hash that the repository records for it. For servers
without version control, the release number was recorded together with the
digests published by the server or recorded by earlier project contracts. Two
dates are kept apart: the date on which the project obtained each file, and the
date at which the source reports that its copy was last changed, such as the
commit date or the last-modification date declared by the server. For example,
the CNNProm benchmark obtained in 2026 corresponds to a repository state from
2016. Any change of source, version, or expected content defines a new
collection version rather than modifying the existing one.

\paragraph{Canonical representation.}
The acquired files were converted into one uniform tabular representation per
released dataset. The conversion was deliberately conservative. It removed
formatting whitespace and normalized letter case, and it neither discarded,
trimmed, deduplicated, nor relabelled any record. Properties that later stages
must handle, namely length deviations, symbols outside $\Sigma_{\mathrm{DNA}}$, exact
duplicates, and identical sequences carrying opposite labels, were recorded as
attributes of each record instead of being corrected silently. A released
dataset was accepted only when its class counts equalled those of the release
and, where a previously used copy existed, when both copies contained the same
multiset of sequences in each class.

\paragraph{Continuity with earlier experiments.}
The experiments reported in this thesis were prepared from copies obtained
before this procedure existed, for which no acquisition record survives. Their
acquisition date is therefore reported as unknown, and it is not inferred from
file-system metadata. Their origin was instead established by evidence. The
preparation pipeline was executed again on the newly acquired files, under the
original rules and random seed. It reproduced, bit for bit, the model-ready
training and test partitions of all eighteen confirmatory datasets, and of
seven of the eight GraphPro datasets, which are retained only for exploratory
use. In the remaining GraphPro dataset, a correction of the organism name
changed the seed from which its partition is derived; the same examples were
recovered, assigned differently between partitions. The confirmatory results
are thus linked to identified upstream releases without being recomputed.

\paragraph{Observed properties of the releases.}
The renewed acquisition exposed release characteristics that are relevant for
interpreting published benchmarks. In one multi-species prokaryotic benchmark,
identical sequences occur with both labels, and some of them span the
distributed training and test partitions. In two eukaryotic promoter
collections, a large fraction of the positive examples are exact duplicates.
One file distributed as a plant promoter set is a mouse promoter set; this case
is examined below. For one bacterial resource, the published window around the transcription
start site does not match the length of the distributed sequences. These
observations are reported as properties of the releases. They motivate the
explicit treatment of duplicates and label collisions summarized in
Table~\ref{tab:sequence-integrity-policies}, and they qualify any comparison with
results reported on the original files.

\paragraph{Identity of a released collection: the Bhandari plant set.}
\citet{bhandari2021comparison} compare learning methods on promoters of yeast,
human, and \textit{Arabidopsis thaliana}, and report results for each organism
and for transfer between them. Their publication states that about 41{,}671
\textit{A. thaliana} promoter regions, from 700 bases upstream to 300 bases
downstream of the transcription start site, were extracted from the UCSC
assembly \texttt{araTha1}, and points to a data folder in which the plant file
is named \texttt{Plant\_promoters}. A name and a declaration do not establish
the organism of a file. The identity of this file was therefore examined with
evidence independent of its name, in four steps: the assembly and chromosome
names declared in its record headers were inspected; sequences were recreated
from the declared assembly at the declared coordinates; the position assumed
to be the transcription start site was compared with the start of the named
transcript in dated annotation releases; and the history of the authors'
repository was searched for earlier releases.

The distributed file does not contain plant sequences. All 106{,}520 headers
name the mouse assembly \texttt{mm10}, the records lie on mouse chromosomes
that have no counterpart in the five-chromosome \textit{A. thaliana} genome,
all 106{,}520 sequences were recreated exactly from the mouse genome at their
declared coordinates, and for
46{,}817 records the assumed start site coincides with the start of the named
mouse transcript. The repository history, in turn, contains a file named
\texttt{Promoter\_ArabidopsisThaliana.csv}, available from April 2020 and
removed in January 2021, on the date of the archives in the data folder that
replaced the repository as distribution point. It holds
exactly 41{,}671 sequences of 1{,}000 bases, the number stated in the
publication, and 41{,}594 of them occur exactly, over their full length, in the
\textit{A. thaliana} reference genome; the human and yeast files removed on the
same day are identical in content to those of the data folder. The plant set
described by the publication therefore exists, but under another name and in
an earlier release than the one the publication points to, while the
designated plant file is a mouse set.

These results determine how each file can be used, following the conditions
defined in Section~\ref{sec:protocol-data-conditions}. The mouse file keeps its
original designation and is used, if at all, as a mouse dataset, with its
verified organism recorded beside the designation; it cannot support any
statement about plants. The earlier file can be used to repeat the published
procedure on the data the publication describes. Such a repetition remains
qualified: the random seed and the shuffled negatives are not available, and
the released notebook selects the first records of each class where the text
describes a random selection. Removing identical sequences, as the notebook does,
also leaves overlapping windows: 13{,}888 of the 37{,}881 distinct located
windows share at least 500 bases with another window, mostly alternative start
sites of one gene, so a random split places overlapping regions in both
training and test sets. Whether the published plant results were
computed on this earlier file is not established, because no output of a plant
run was released. In addition, the 77 sequences not found in the current
reference all occur exactly in the earlier mitochondrial genome sequence, which
the current reference replaced in 2018, so every sequence of the file is
located in an \textit{A. thaliana} reference.

Whether the windows are centred on transcription start sites depends on the
annotation used to judge them. Of the 41{,}591 windows that have the declared
length and a genome locus, 22{,}970 place the declared start site exactly at
the start of a transcript in the current RefSeq annotation, and similar
numbers hold for its 2018 and 2023 releases. Most of the remaining windows lie
within a few hundred bases of an annotated start, and the annotated start is
far more often upstream than downstream of the window. Against the TAIR10 gene
models, however, which contain exactly 41{,}671 transcripts, the declared start
site coincides with a transcript start in 41{,}539 windows; the other 52 lie on
the mitochondrial genome or on a nuclear segment of mitochondrial origin. The
apparent disagreements therefore reflect later revisions of the annotation,
which extended many transcripts at their 5$'$ end, rather than errors in the
windows, although this interpretation, and the correspondence of each window
to one transcript, have not been tested record by record. For the released
condition, the file was converted into a canonical table that keeps every
record, sequence, and label, with an identifier derived from its row in the
acquired file; writing the table back reproduced the acquired file byte for
byte. Genome coordinates and start-site classes are kept apart from that table,
as evidence for a possible reconstruction, and do not alter it.

Repeating a published computation and relying on its conclusions are different
claims. A repetition on the identified plant file can show that the reported
values are reproducible; it cannot remove the limitations of the original
design, such as synthetic negatives and duplicate sequences, which affect what
those values mean. Whether the identity error changes the conclusions of model
comparisons is an empirical question addressed by evaluating the same models
under the released and reconstructed conditions.

\paragraph{Limitations.}
For servers that publish no checksums, the identity of a file rests on its first
acquisition by the project, supplemented by agreement with earlier copies where
these exist. Several servers report no modification date, so the version of
their content can be dated only by acquisition. Exact agreement of sequences
does not establish their biological annotation: the procedure guarantees that
the data are the released data, not that the released labels are correct.

